% !TEX root = ../termpaper.tex
% @author Leonhard Lüdeke
%

\section{Einleitung}

\subsection{Das Spiel Vier Gewinnt}
Vier Gewinnt ist ein klassisches Strategiespiel, und wurde 1973 von Howard Wexler handelt sich um ein endliches Zwei-Personen Nullsummenspiel mit perfekter Information. Ein Nullsummenspiel kennzeichnet sich dadurch aus, dass die Verluste des einen Spielers exakt den Gewinnen des anderen darstellen. Eine spielende Person ist deshalb perfekt Informiert, da die gesamte Spielsituation jederzeit offen einsehbar ist. Somit kann eine spielende Person eine voll informierte Entscheidung treffen. \cite{owen2013spieltheorie}
Das Spiel wird auf einem vertikal aufgerichteten Brett mit sieben Spalten und sechs Zeilen gespielt. Die Spielenden werfen abwechselnd Steine ihrer Farbe von oben in eine Spalte ein, wobei der Stein auf die unterste freie Position fällt. Gewinner ist derjenige Person, die zuerst vier eigene Steine horizontal, vertikal oder diagonal in einer Linie anordnet. \cite{wiki:Vier_gewinnt} 

Die Kombination aus einfacher Spielmechanik und überraschend hohen Gesamtzahl der validen Spielpositionen von circa 4,5 Billionen \cite{dengel2008ki}, während Tic-Tac-Toe 5.478 gültige Positionen besitzt. Daher ist ein Aufbau eines Kompletten Spielbaumes zwar nicht unmöglich, jedoch ist wäre dieser Baum mit $4 x 10^{21}$ Blattknoten, so groß das selbst im Komprimierten zustand 557 Gigabyte nötig wären um diesen Spielbaum abzubilden \cite{yokota2022high}. Zusätzlich dürfte die Laufzeit für die Berechnung des Spielbaums viel Zeit in Anspruch nehmen. Die Lösung des Spiel mittels eines Gamebot's liegt somit nicht in einer 1 zu 1 Abbildung aller Spielzustände, sondern in einem Algorithmischen Ansatz, welche ich im nächsten Unterkapitel beleuchtet werden.


\subsection{Algorithmische Lösungsansätze für einen Bot}
Wie oben beschrieben gebe es die Möglichkeit einen Brute-Force-Ansatz zu verfolgen, jedoch würde das einen exponentiellen Wachstum des Suchraumes bedeuten und schlichtweg zu unpraktikabel sein. Jedoch gibt es einen Algorithmus der zur Berechnung optimaler Strategien bei Nullsummenspielen. Dies machte sich Victor Allis in der Masterarbeit "A Knowledge-based Approach of Connect-Four" \cite{allen1990expert} zunutze und zeigte in dieser Arbeit das die Beginnende Spielende Person immer gewinnt, sofern sie keinen Fehler macht. Zeitgleich gelang diese auch James D. Allen \cite{allen1990expert} 1998.

 


Wie James D. Allen \cite{allen1990expert} und Victor Allis im Jahre 1988 unabhängig voneinander zeigten, gibt es eine Möglichkeit für den Beginnenden Spieler immer zu gewinnen. Victor Allis nutzte in der Masterarbeit „A Knowledge-based Approach of Connect-Four" \cite{allis1988knowledge} einen Minimax-Algortihmus mit spezifischen Heuristiken. Dies klingt recht vielversprechende



Mit Hilfe dieser Methoden lassen sich die Grenzen klassischer, deterministischer Verfahren aufzeigen und mit stochastischen oder heuristischen Ansätzen vergleichen. Ziel dieser Arbeit ist es, die jeweiligen Stärken und Schwächen der Ansätze im Kontext von Vier Gewinnt aufzuzeigen und praktische Einsatzmöglichkeiten zu evaluieren.
